Tento výzkumný úkol se zabývá teorií a návrhem propojení umělé inteligence (AI) a strojového učení ve hrách k dynamickému přizpůsobování chování protivníků v reálném čase.

V první kapitole jsou obsaženy dvě části. První část popisuje umělou inteligenci ve hrách. Specificky popisuje historii umělé inteligence ve stolních a digitálních hrách, dále rozebírá některé algoritmy, které jsou nejčastěji používané k vytvoření umělé inteligence v digitálních hrách, ale díky kterým není umělá inteligence schopna se přizpůsobit nebo upravit svoje chování. Nakonec se tato část specificky zaměří na některé existující 3D akční hry, které obsahují inteligentnější AI.
Druhá část obsahuje základy strojového učení, existující typy strojového učení a také některé algoritmy z těchto typů. 

%Kapitola druhá také obsahuje dvě podkapitoly. První podkapitola obsahuje existující řešení na tuto problematiku v~herním průmyslu, kde první ze zmíňených řešení je autonomní agent postaven tak, aby za vás dokázal hrát danou hru. Druhé řešení je systém ve hře jménem Shadow of War, kde tento systém běží na pozadí, když je potřeba, a sbírá data z~bitev, aby upravil vzhled, vztah nebo atributy nepřítele a znova ho vyslal zaútočit na hlavní postavu. Ve druhé podkapitole je zahrnut návrh umělé inteligence a vizualizace a také použitý software a knihovny potřebný k~vytvoření vizualizace a umělé inteligence.

%Závěr popisuje úspěšnost bakalářské práce, samotnou hru a také případné budoucí rozšíření, pomocí kterého by se dal autonomní agent vylepšit.